\Subhead{1.3.3}{Применение захватывания для стабилизации частоты}
Маломощный стабильный источник может захватить мощный, но менее стабильный генератор. В установившемся захвате частота выхода равна $\Omega$; медленный уход $\omega_g$ компенсируется изменением $\theta_*$. Условие --- дрейф остаётся внутри полосы захвата.

При малом изменении собственной частоты линеаризация \eqref{eq:adler} даёт
\begin{equation}
 \delta\dot\theta+K\cos\theta_*\delta\theta=\delta\omega_g(t).
\end{equation}
Она показывает время слежения: быстрые флуктуации подавляются иначе, чем медленные, а шум эталона переносится на выход. На границе полосы возможны проскальзывания и срыв. Примеры --- СВЧ-генератор с внешним захватом и ведущий/ведомый лазеры; в последнем случае связываются амплитуды \emph{поля}, а не непосредственно мощности.
